\chapter*{The Argument of Part IV}
\pdfbookmark[0]{The Argument of Part IV}{part4.guide}
\label{part4:guide}
\markboth{Open Source Evaluatology}{The Argument of Part IV}
\noindent The preceding parts established the conditions under which open source observations become evaluative evidence. Part I specified the question, value, and research design. Part II developed Activity, OpenRank, and OpenDigger. Part III examined recognition, continuity, and governance in settings where a judgment matters to participants. The final part asks how this knowledge can support institutions and sustain a field of inquiry.

Chapter 10 follows the work of an open source science and technology think tank. A decision question determines the relevant observation; a report or interface makes that observation interpretable; an institutional record establishes whether and how it was used. The chapter connects the OpenThink architecture with global and regional analyses, community applications, and OpenDipper. Its central contribution is an account of how evidence travels into a consequential setting while retaining its meaning.

Chapter 11 develops the conditions for making this work cumulative. The six disciplinary capabilities organize theories, data, methods, benchmarks, tools, and evaluation activities. The tool studies expose different ways of working with related evidence. OpenPerf supplies concrete comparisons, and the interdisciplinary synthesis relates the book's findings to measurement, networks, software work, education, and institutions. These connections support a research programme grounded in the studies already presented.

\vspace{5pt}
\begingroup\small
\renewcommand{\arraystretch}{1.17}
\noindent\begin{tabularx}{\linewidth}{@{}>{\raggedright\arraybackslash}p{29mm}>{\raggedright\arraybackslash}X >{\raggedright\arraybackslash}X@{}}
\toprule
\textbf{Connection} & \textbf{Chapter 10} & \textbf{Chapter 11}\\
\midrule
Question and evidence & Select observations for a specified institutional decision. & Preserve the evaluation question and its methodological commitments across studies.\\[5pt]
Tools and products & Use reports, maps, dashboards, and project investigations to communicate findings. & Examine the roles, semantics, and exchange of evidence across an open tool ecosystem.\\[5pt]
Learning from use & Distinguish dissemination, uptake, and consequences; retain the record of interpretation. & Connect those requirements to validation, shared benchmarks, and theoretical revision.\\
\bottomrule
\end{tabularx}
\endgroup
\par\vspace{9pt}
The two chapters operate at different levels. A think-tank product answers a situated question; a disciplinary framework organizes the knowledge and practices that make such answers possible. The five-layer OpenThink architecture therefore retains its institutional purpose, while Chapter 11's six-capability framework describes the field. Both draw on the five-component evaluation-system design in Chapter 3.

The empirical record remains versioned. Report-year observations are not merged with calendar-year snapshots, OpenGalaxy's 2024 research platform is distinguished from OpenDipper's later Open Galaxy module, and OpenPerf's published experiments retain their original conditions. These distinctions make the final synthesis possible: the book connects results through their questions and implications while preserving the evidence on which each rests.
\clearpage
